\section{The remaining higher-block reach question}\label{sec:higher-reach}

No result in the preceding sections assumes an exact reach theorem beyond
four blocks. We now give the precise unresolved route and an exact test
of one proposed relaxation. This separates a failed implication between
linear inequalities from a counterexample generated by an actual control.

\subsection{An exclusion identity and its missing premise}

Fix $n\ge k+1$, $k\ge3$, $E>0$, and reach maps on $[0,L]$.
For an available set $U$ and $h\le |U|$, let $M_U^h$ be the maximum
reach over all ordered $h$-tuples of distinct elements of $U$.
Write $M^h=M_{[n]}^h$, $M_i^h=M_{[n]\setminus\{i\}}^h$, and
$M_{ij}^h=M_{[n]\setminus\{i,j\}}^h$. These are actual maxima of
reach compositions, not independent event variables. Assume
$M^{k-1}<L$ so every event used below is uncapped.

\begin{lemma}[General exclusion inequality]\label{lem:general-exclusion}
Let $S$ be a $(k-1)$-element set supporting a word attaining $M^{k-1}$.
Then
\begin{align}
 (n-2)M_i^{k-1}
 &\ge nE+M_i^{k-2}+\sum_{j\ne i}M_{ij}^{k-2}-M^{k-1},
       \label{eq:general-exclusion}\\
 \sum_iM_i^{k-1}
 &\ge\left(n-k+1-\frac{k-1}{n-2}\right)M^{k-1}
       +\frac{(k-1)nE+\mathcal H_S^{k-2}}{n-2},
       \label{eq:general-exclusion-sum}\\
 \mathcal H_S^{k-2}
 &:=\sum_{i\in S}\left(M_i^{k-2}+\sum_{j\ne i}M_{ij}^{k-2}\right).
       \label{eq:general-weighted-H}
\end{align}
If also $M^{k-1}\ge B_{k-1}$ and
$\mathcal H_S^{k-2}\ge(k-1)nB_{k-2}$, then $k$ distinct blocks reach
$L$ whenever $L\le B_k$.
\end{lemma}
\begin{proof}
Appending $j$ to a maximizing $(k-2)$-word avoiding $i,j$ gives a
reach no larger than $M_i^{k-1}$. By the uncapped boundary identity
and monotonicity,
\[
 A_j(M_i^{k-1})\le M_i^{k-1}-E-M_{ij}^{k-2}\qquad(j\ne i).
\]
Sum over $j\ne i$ and subtract from conservation to obtain
\[
 A_i(M_i^{k-1})\ge(n-1)E+\sum_{j\ne i}M_{ij}^{k-2}
                              -(n-2)M_i^{k-1}.
\]
Appending $i$ to a maximizing word counted by $M_i^{k-2}$ likewise
gives $A_i(M^{k-1})\le M^{k-1}-E-M_i^{k-2}$.
Since $M_i^{k-1}\le M^{k-1}$, monotonicity of $A_i$ combines these
bounds into~\eqref{eq:general-exclusion}. For $i\notin S$ the same
maximizing word remains available, so $M_i^{k-1}=M^{k-1}$.
Sum~\eqref{eq:general-exclusion} over $i\in S$ and add these equal
outside terms to obtain~\eqref{eq:general-exclusion-sum}.

Its coefficient of $M^{k-1}$ is positive for $n\ge k+1$.
Under the additional premises, use
$n(E+B_{k-2})=(n-1)B_{k-1}$ to deduce
$\sum_iM_i^{k-1}\ge nB_{k-1}$. If no $k$-word reaches $L$, append
$i$ to a maximizing word avoiding $i$. Failure implies the strict
inequality $H_i(L)>M_i^{k-1}+E$. Summing gives
\[
 (n-1)L>\sum_iM_i^{k-1}+nE
             \ge nB_{k-1}+nE=(n-1)B_k,
\]
a contradiction to $L\le B_k$.
\end{proof}

For $k=5$, the premise $M^4\ge B_4$ is already established whenever
$M^4<L$. The missing premise in this route is the weighted three-block
inequality $\mathcal H_S^3\ge4nB_3$ for a maximizing four-mode set
$S$. The stronger statement for every four-element $S$ would suffice
as well. For larger $k$ both the corresponding weighted inequality and
the preceding exact reach induction need justification.
The research question is whether every $n>k\ge5$ measurable input admits
$k$ distinct blocks reaching $\min\{T,B_k(n,E)\}$ for every $E>0$.
If true, it would identify $G^-_{n,k}(T)=TL_{n,k}$ and, by the proved
heavy-mode reduction, the full minimax. The question is unused, and the
boundary example at $n=k=3$ does not answer it in the spare-mode range.

\subsection{A nonphysical feasible event assignment}\label{subsec:negative-relaxation}

Take $n=6$, $E=1$, label modes $0,\ldots,5$, and fix $S=\{0,1,2,3\}$.
A natural relaxation stores six roots $R_i$ and, for every available set
$U$ of size at least three, a pair event $(2,U)$; for every $|U|\ge4$
it stores a triple event $(3,U)$. Each event $v$ has time $t_v$ and
six allocations $a_{v,i}$. There are 42 pair events and 22 triple events,
hence $6+7\cdot64=454$ nonnegative variables. The full row specification is
\begin{align*}
 &R_i\ge1,\quad R_0\le R_i\ (i\ne0),\quad\sum_{i\ne0}R_i\ge6,\\
 &\sum_i a_{v,i}=t_v,\quad a_{(h,U),i}\ge R_i-1\quad(i\in U),\\
 &t_{(2,U)}-a_{(2,U),i}\ge R_j+1\quad(i,j\in U,\ i\ne j),\\
 &t_{(3,U)}-a_{(3,U),i}\ge t_{(2,U\setminus\{i\})}+1\quad(i\in U),\\
 &a_{(h,U),i}\le a_{(h',U'),i}
       \quad(h\le h',\ U\subseteq U',\ (h,U)\ne(h',U')),\\
 &a_{(h,[6]),i}\le a_{(h,U),i}
       \quad(P_h\subseteq U\ne[6]),
       \qquad P_2=\{0,1\},\ P_3=\{0,1,2\}.
\end{align*}
Here $[6]$ denotes all six labels in this subsection. The last rows and
the inclusion rows together force equality for events sharing the fixed
global maximizing pair or triple. Each row is necessary for actual
uncapped events with those maximizers and with $R_0$ smallest. Indeed,
adding available modes or activation blocks cannot decrease maximum reach,
and coordinate allocations increase with time. In the all-available
three-block objective define
\[
 \mathcal H_S^3=\sum_{i\in S}\left(t_{(3,[6]\setminus\{i\})}
               +\sum_{j\ne i}t_{(3,[6]\setminus\{i,j\})}\right).
\]
The bundled rational assignment satisfies all 3660 displayed inequalities
and all 64 conservation equalities, with nonnegative variables, but gives
\begin{equation}\label{eq:negative-event-objective}
 \mathcal H_S^3=\frac{40328}{387}<\frac{13104}{125}=4nB_3(6,1).
\end{equation}
Both a row-matrix checker and an independently organized event checker
verify these facts exactly. Thus these necessary constraints do not imply
the proposed weighted inequality for arbitrary four-element $S$.

This is not a counterexample to any reach theorem. The pair event for
$U=\{0,2,3\}$ has time $1658/645$; that for $V=\{2,3,4,5\}$ has
the later time $614/215$. Yet the corresponding mode 2 allocation falls
from $239/645$ to $47/129$. The time difference is $184/645>0$ and
the allocation drop is $4/645>0$. The sets are incomparable by inclusion,
so the relaxation omitted this chronological comparison. A common
cumulative control cannot produce these data. Nor does the relaxation
require $S$ itself to support a global maximizing four-word, a further
condition if one only seeks the weaker premise in
Lemma~\ref{lem:general-exclusion}.

\begin{lemma}[Interpolation of event allocations]\label{lem:event-interpolation}
Let finitely many events $(t_v,a_v)$ have $t_v\ge0$, $a_v\in\R_+^n$,
and $\sum_i a_{v,i}=t_v$. They lie on a common admissible cumulative
trajectory if and only if
\begin{equation}\label{eq:chronological-order}
 t_v\le t_w\quad\Longrightarrow\quad a_{v,i}\le a_{w,i}
                    \quad\hbox{for every }i.
\end{equation}
Such a trajectory can be chosen piecewise affine, up to and beyond the
last event.
\end{lemma}
\begin{proof}
Necessity follows from coordinate monotonicity. For sufficiency, sort the
events by time and add $(0,0)$. Equal-time allocations agree because
coordinate order in both directions forces equality; equivalently,
nonnegative ordered differences with zero total are zero. Between distinct
successive times interpolate linearly. Each coordinate increment is
nonnegative and their sum is the time increment, so the slope vector is
simplex-valued. Extend after the last event with any simplex-valued constant
rate. This gives an admissible trajectory containing all event data.
\end{proof}

This lemma characterizes trajectory interpolation only. The variables
$R_i,t_{h,U}$ must additionally equal the specified latest inverse and
maximum-composition reaches of that same trajectory. Conservation and
chronological interpolation alone do not establish these identities.

\subsection{An exact chronological-chamber test}\label{subsec:chronological-test}

We strengthened the failed relaxation by sorting its 64 events in the
nondecreasing time order of the saved witness. Ties are resolved by their
original enumeration: first order $h=2,3$, then decreasing $|U|$, then
lexicographic order of the increasing mode tuples. Let $v_1,\ldots,v_{64}$
be that fixed permutation. Add the 378 inequalities
\begin{equation}\label{eq:chamber-cuts}
 a_{v_j,i}\le a_{v_{j+1},i}\qquad(1\le j<64,\ 0\le i<6).
\end{equation}
Conservation then also forces $t_{v_j}\le t_{v_{j+1}}$.
Thus this is one closed chronological chamber, allowing ties, not a
global imposition of the old witness's numerical times. The witness itself
has 239 violations of ordered coordinate comparisons, with largest drop
$224/645$, and is excluded by these new rows.

\begin{proposition}[One strengthened chamber, computer-assisted]
\label{prop:chronological-chamber}
For this single fixed event permutation, the strengthened relaxation with
454 variables, 4038 inequalities, and 64 equalities has exact minimum
\[
 \min\mathcal H_S^3=13104/125=4nB_3(6,1).
\]
\end{proposition}
\begin{proof}
Write the inequalities as $Ax\le b$, the equalities as $Qx=d$, and
the objective as $c^\top x$, with $x\ge0$. The supplement supplies
rational vectors $y,z$ and verifies, component by component,
\[
 y\le0,\qquad c-A^\top y-Q^\top z\ge0,\qquad
 b^\top y+d^\top z=13104/125.
\]
Consequently every feasible $x$ satisfies
$c^\top x\ge y^\top Ax+z^\top Qx\ge b^\top y+d^\top z$.
There are 254 nonzero dual rows; the checker reconstructs every integer
row from the bundled original relaxation and the explicit permutation,
then sorts normalized sparse rows and their right-hand sides
lexicographically before applying the dual vectors.
It uses only rational arithmetic and does not invoke a numerical solver.

For attainment, take $R_i=6/5$, every pair-event time $66/25$, every
triple-event time $546/125$, and $a_{v,i}=t_v/6$. These are precisely the
uniform-control events. In the chosen permutation all pair events precede
all triple events, so every chamber row holds. All original rows and the
objective equality are verified exactly as well. There are 24 triple-time
terms in the objective, giving $24(546/125)=13104/125$.
\end{proof}

The exact certificate closes the relaxation gap in this chamber. It does
not cover the other chronological permutations, other choices of global
maximizers, or other mode counts. It also does not claim that every point
in the strengthened relaxation satisfies the inverse reach identities.
A floating-point linear program was used only to discover the rational
dual; the proof is the checked rational identity and the matching explicit
primal. The general five-block reach problem remains unresolved by this
investigation. Separately, the possible strengthening of
Lemma~\ref{lem:first-repeat} to force an adjacent pair of a largest heavy
mode remains unused and unproved; it is independent of the established
existential heavy-mode theorem.
